\documentclass[a4paper,12pt]{article}
\usepackage[utf8]{inputenc}
\usepackage[T1]{fontenc}
\usepackage[ngerman]{babel}
\usepackage{amsmath, amssymb}
\usepackage{geometry}
\usepackage{tikz}
\usetikzlibrary{shapes.geometric, arrows.meta, positioning}
\usepackage{parskip}
\usepackage{titlesec}
\usepackage{enumitem}

\geometry{left=3cm, right=3cm, top=2.5cm, bottom=2.5cm}
\titleformat{\section}{\large\bfseries}{}{0em}{}[\titlerule]
\titleformat{\subsection}{\normalsize\bfseries}{}{0em}{}

\tikzset{
  box/.style={rectangle, draw, rounded corners=3pt, text width=4.8cm,
              align=center, minimum height=1.0cm, font=\small, inner sep=4pt},
  arrow/.style={-{Latex[length=2mm]}, thick},
}

\begin{document}

\begin{center}
  {\LARGE\bfseries Studierhinweise zu Lektion 5}\\[0.5em]
  {\large Lineare Algebra (Modul 61112)}
\end{center}

\vspace{0.8em}

Diese Lektion stellt die Werkzeuge für das \textbf{Normalformproblem} bereit: Zu einem Endomorphismus
$\varphi \in \mathrm{End}(V)$ sucht man eine Basis, in der die Matrixdarstellung möglichst einfach
(„normal") ist. Der zentrale Begriff ist der des \textbf{Eigenvektors} bzw.\ \textbf{Eigenwertes}.
Eigenwerte lassen sich als Nullstellen des \textbf{charakteristischen Polynoms} berechnen.
Das \textbf{Minimalpolynom} und der \textbf{Satz von Cayley-Hamilton} runden die Theorie ab.

Die Begriffe in Abschnitt~5.2 sind fundamental -- nehmen Sie sich dafür ausreichend Zeit,
da sie in Lektion~6 ständig benötigt werden.

% -----------------------------------------------------------------------
\section*{Struktur der Lektion 5}

\begin{center}
\begin{tikzpicture}[node distance=0.55cm]
  \node[box] (nfp)  {5.1 Normalformproblem\\Ähnlichkeit, invariante Unterräume};
  \node[box, below=0.5cm of nfp] (ev)  {5.2 Eigenwerte, Eigenvektoren\\Eigenräume, Diagonalisierbarkeit};
  \node[box, below=0.5cm of ev]  (pol) {5.3 Char.\ Polynom, Minimalpolynom\\Satz von Cayley-Hamilton};
  \node[box, below=0.5cm of pol] (sto) {5.4 Stochastische Matrizen\\Satz von Perron, PageRank};
  \draw[arrow] (nfp) -- (ev);
  \draw[arrow] (ev) -- (pol);
  \draw[arrow] (pol) -- (sto);
\end{tikzpicture}
\end{center}

\newpage

% -----------------------------------------------------------------------
\section*{Zielelement 5.1 -- Normalformproblem und invariante Unterräume}

\subsection*{Lerninhalte}
\begin{center}
\begin{tikzpicture}[node distance=0.55cm]
  \node[box] (äh)  {Normalformproblem:\\Ähnlichkeit von Matrizen};
  \node[box, below=0.5cm of äh]  (inv) {Invarianter Unterraum $U$:\\$\varphi(U) \subseteq U$, Einschränkung $\varphi|_U$};
  \node[box, below=0.5cm of inv] (zer) {Zerlegung in invariante Unterräume\\$\Rightarrow$ blockdiagonale Matrix (5.1.13)};
  \draw[arrow] (äh) -- (inv);
  \draw[arrow] (inv) -- (zer);
\end{tikzpicture}
\end{center}

\subsection*{Lernziele}
\begin{itemize}
  \item Das Normalformproblem für Endomorphismen und Matrizen formulieren.
  \item Den Begriff des $\varphi$-invarianten Unterraums kennen und Beispiele angeben.
  \item Verstehen, warum eine Zerlegung in invariante Unterräume zu einer blockdiagonalen Matrixdarstellung führt (Satz~5.1.13).
\end{itemize}

\subsection*{Selbstkontrollelement 5.1}
Sei $A = \begin{pmatrix}2&1\\0&3\end{pmatrix}$. Zeigen Sie, dass $\langle e_1 \rangle$ ein $A$-invarianter Unterraum ist.

\newpage

% -----------------------------------------------------------------------
\section*{Zielelement 5.2 -- Eigenwerte, Eigenvektoren, Diagonalisierbarkeit}

\subsection*{Lerninhalte}
\begin{center}
\begin{tikzpicture}[node distance=0.55cm]
  \node[box] (def) {Eigenvektor $v \neq 0$, Eigenwert $\lambda$:\\$\varphi(v) = \lambda v$};
  \node[box, below=0.5cm of def] (era) {Eigenraum $E_\lambda(\varphi) = \ker(\varphi - \lambda \mathrm{id})$\\Geometrische Vielfachheit};
  \node[box, below=0.5cm of era] (lem) {Lemma 5.2.13: Summe der Eigenräume\\ist direkt};
  \node[box, below=0.5cm of lem] (dia) {Diagonalisierbar $\iff$ Basis aus Eigenvektoren};
  \draw[arrow] (def) -- (era);
  \draw[arrow] (era) -- (lem);
  \draw[arrow] (lem) -- (dia);
\end{tikzpicture}
\end{center}

\subsection*{Lernziele}
\begin{itemize}
  \item Eigenwerte und Eigenvektoren definieren; Eigenräume berechnen.
  \item Die geometrische Vielfachheit eines Eigenwertes bestimmen.
  \item Das Lemma~5.2.13 kennen: Eigenräume zu verschiedenen Eigenwerten stehen in direkter Summe.
  \item Kriterium für Diagonalisierbarkeit kennen und anwenden.
\end{itemize}

\subsection*{Selbstkontrollelement 5.2}
Bestimmen Sie alle Eigenwerte und Eigenräume von $A = \begin{pmatrix}3&1\\0&2\end{pmatrix}$.
Ist $A$ diagonalisierbar?

\newpage

% -----------------------------------------------------------------------
\section*{Zielelement 5.3 -- Charakteristisches Polynom und Minimalpolynom}

\subsection*{Lerninhalte}
\begin{center}
\begin{tikzpicture}[node distance=0.55cm]
  \node[box] (cp)  {Char.\ Polynom $\chi_\varphi(X) = \det(X \cdot \mathrm{id} - \varphi)$\\Nullstellen = Eigenwerte};
  \node[box, below=0.5cm of cp]  (av)  {Algebraische Vielfachheit\\eines Eigenwertes};
  \node[box, below=0.5cm of av]  (mp)  {Minimalpolynom $\mu_\varphi$:\\kleinster normierter Teiler von $\chi_\varphi$\\mit $\mu_\varphi(\varphi) = 0$};
  \node[box, below=0.5cm of mp]  (ch)  {Satz von Cayley-Hamilton (5.3.42):\\$\chi_\varphi(\varphi) = 0$, d.\,h.\ $\mu_\varphi \mid \chi_\varphi$};
  \draw[arrow] (cp) -- (av);
  \draw[arrow] (av) -- (mp);
  \draw[arrow] (mp) -- (ch);
\end{tikzpicture}
\end{center}

\subsection*{Lernziele}
\begin{itemize}
  \item Das charakteristische Polynom $\chi_\varphi$ berechnen; Grad ist $\dim V$.
  \item Algebraische und geometrische Vielfachheit unterscheiden; $1 \le \text{geom.} \le \text{alg.}$
  \item Das Minimalpolynom als normierten Erzeuger des Verschwindungsideals (5.3.32) kennen.
  \item Den Satz von Cayley-Hamilton (5.3.42) anwenden: $\chi_\varphi(\varphi) = 0$.
\end{itemize}

\subsection*{Selbstkontrollelement 5.3}
Berechnen Sie $\chi_A$ und $\mu_A$ für $A = \begin{pmatrix}2&1\\0&2\end{pmatrix}$.
Verifizieren Sie den Satz von Cayley-Hamilton.

\newpage

% -----------------------------------------------------------------------
\section*{Zielelement 5.4 -- Stochastische Matrizen und PageRank}

\subsection*{Lerninhalte}
\begin{center}
\begin{tikzpicture}[node distance=0.55cm]
  \node[box] (sto) {Stochastische Matrix:\\Spalten summieren sich zu 1,\\alle Einträge $\geq 0$};
  \node[box, below=0.5cm of sto] (per) {Satz von Perron (5.4.6):\\Eigenwert $1$ mit geo.\ Vfh. $1$};
  \node[box, below=0.5cm of per] (pag) {Google PageRank:\\Stationäre Verteilung als Eigenvektor};
  \draw[arrow] (sto) -- (per);
  \draw[arrow] (per) -- (pag);
\end{tikzpicture}
\end{center}

\subsection*{Lernziele}
\begin{itemize}
  \item Stochastische Matrizen definieren und erkennen.
  \item Den Satz von Perron (5.4.6) kennen: positive stochastische Matrizen haben Eigenwert~$1$ mit geometrischer Vielfachheit~$1$.
  \item Das Prinzip des Google PageRank-Algorithmus erklären.
\end{itemize}

\subsection*{Selbstkontrollelement 5.4}
Bestimmen Sie den Eigenvektor zum Eigenwert~$1$ von $A = \begin{pmatrix}0{,}7 & 0{,}4 \\ 0{,}3 & 0{,}6\end{pmatrix}$.

\end{document}
